\section{Methodology}
\subsection{Research questions}
\begin{description}
 \item[RQ1:] Which selected-system faults escape generic structural verification?
 \item[RQ2:] What does due-boundary obligation discharge add beyond invalidation records and lazy demand-driven recomputation?
 \item[RQ3:] How do selective and always verification compare in detection, localization, false positives, and verifier work on the evaluated corpora?
 \item[RQ4:] Does the policy distinction survive source-to-result execution and a second public-SDK language package?
\end{description}

\subsection{Frozen boundary corpus}
The historical System W corpus is immutable: 40 instances representing 32 primary operator shapes in five families, 10 post-freeze challenge operators, and 100 valid controls per policy. Faults are repeated three times. The v3 study preserves historical B0/B1/B2 labels, emits a separate four-policy schema, and mechanically checks the frozen mutation catalog.

\subsection{Demand-driven witness}
To compare against a stronger invalidation baseline without changing the frozen matrix, we execute one controlled invalidation through three schedules: lazy demand with no consumer request, lazy demand with an explicit request, and P2 obligation discharge. The same AIR verifier and invalid artifact are used in the demanded and P2 cells, plus a matched valid control. The witness tests the semantic distinction, not population prevalence.

\subsection{System W source-to-result corpus}
The end-to-end corpus contains 30 programs across arithmetic, pricing, and backend strata. Five cases enable an optimizer that preserves structural validity while violating a declared semantic relation at optimized AIR. Each case/policy pair runs twice in a fresh process, yielding 240 raw records. CIL faults produce a wrong result under P0/P1; interpreter faults may reach a later backend/runtime failure. P2/P3 must reject at optimized AIR with the expected capability diagnostic. After repairing mixed numeric promotion, all 25 non-fault programs remain unchanged and must produce their expected result under every policy.

\subsection{TensorRules and ablations}
TensorRules freezes two valid programs, two intrinsically invalid programs, and eight post-verification mutations of dimensions, output shape, layout, dynamic extent, backend capability, broadcast relation, stale shape, and fact ownership. Four policies produce 48 observations. Separately, eight minimal counterexamples and matched controls isolate producer identity, source-node identity, selected pipeline order, canonical ownership, conflict rejection, unresolved obligations, backend capabilities, and repeated-occurrence rejection.

\subsection{Evidence and claim control}
Raw JSONL is written before assertions. Artifacts carry source identity, environment records, and SHA-256 manifests. Hosted CI is a reproducibility provider, not a performance environment. An external-author freeze/import kit exists, but no externally authored corpus is included in the reported results; corresponding generalization claims remain blocked.
